\begin{remark}\label{thm:massey-products}
阶 $d>2$ 的\trComposition{}不是\trChainMap。不过，它们在\trCohomology{}层面上有部分的化身，
即\textbf{\trMasseyProduct}（Massey product）。设有\trMorphism{}%
$[a_k]\in\Hom_{H(\scr{A})}(X_{k-1},X_k)$（$k=1$，$2$，$3$），使得乘积 $[a_3]\cdot[a_2]$
与 $[a_2]\cdot[a_1]$ 都是零。取链代表 $a_k$，并且另取
$h_k\in\hom_{\scr{A}}(X_{k-1},X_{k+1})$（$k=1$，$2$），使得
$\mu^1_{\scr{A}}(h_k)=\mu^2_{\scr{A}}(a_{k+1},a_k)$。那么
$c=\mu^3_{\scr{A}}(a_3,a_2,a_1)-\mu^2_{\scr{A}}(h_2,a_1)-\mu^2_{\scr{A}}(a_3,h_1)
\in\hom_{\scr{A}}(X_0,X_3)$ 是一个\trDegree{}为 $|a_1|+|a_2|+|a_3|-1$ 的\trCocycle，而类
\begin{equation}\label{eq:massey-class}
[c]\in\frac{\Hom_{H(\scr{A})}(X_0,X_3)}
{[a_3]\cdot\Hom_{H(\scr{A})}(X_0,X_2)+\Hom_{H(\scr{A})}(X_1,X_3)\cdot[a_1]}
\end{equation}
只依赖于原先的 $[a_k]$。通常写作
$\langle[a_3],[a_2],[a_1]\rangle=(-1)^{|a_2|}[c]$。
\end{remark}

\paragraph{\trFunctor。}\label{sub:functors}
两个\trNonunital{}\trACategory{} $\scr{A}$、$\scr{B}$ 之间的
\textbf{\trNonunital{}\trAFunctor}（non-unital $A_\infty$-functor）由一个映射
$\scr{F}\colon\Ob\scr{A}\to\Ob\scr{B}$ 与各阶 $d\ge 1$ 的\trMultilinear{}映射组成：
\begin{equation}\label{eq:afunctor-components}
\scr{F}^d\colon\hom_{\scr{A}}(X_{d-1},X_d)\otimes\cdots\otimes\hom_{\scr{A}}(X_0,X_1)
\longrightarrow\hom_{\scr{B}}(\scr{F}X_0,\scr{F}X_d)[1-d]
\end{equation}
其中 $X_0,\ldots,X_d\in\Ob\scr{A}$，要求满足下列（多项式型的）方程：
\begin{equation}\label{eq:afunctor-equations}
\begin{split}
&\sum_r\sum_{s_1,\ldots,s_r}\mu^r_{\scr{B}}
\bigl(\scr{F}^{s_r}(a_d,\ldots,a_{d-s_r+1}),\ldots,\scr{F}^{s_1}(a_{s_1},\ldots,a_1)\bigr)
\\
&\hspace{2em}=\sum_{m,n}(-1)^{\clubsuit_n}\,\scr{F}^{d-m+1}
(a_d,\ldots,a_{n+m+1},\mu^m_{\scr{A}}(a_{n+m},\ldots,a_{n+1}),a_n,\ldots,a_1),
\end{split}
\end{equation}
其中左边的和取遍所有 $r\ge 1$ 与\trPartition{} $s_1+\cdots+s_r=d$。

有一个诱导出的普通（\trLinear、\trGraded、\trNonunital）\trFunctor{}%
$H(\scr{F})\colon H(\scr{A})\to H(\scr{B})$，它在\trMorphism{}上的作用是
$[a]\mapsto[\scr{F}^1(a)]$。若 $H(\scr{F})$ 是\trCategory{}之间的\trIsomorphism，就称
$\scr{F}$ 为\textbf{\trQuasiisomorphism}（quasi-isomorphism）；若 $H(\scr{F})$ 是%
\trFullyFaithful{}的，就称 $\scr{F}$ 为\textbf{\trCohomology{}\trFullyFaithful}
（cohomologically full and faithful）。

\trNonunital{}\trAFunctor{}可以按如下方式\trComposition：
\begin{equation}\label{eq:afunctor-composition}
\begin{split}
&(\scr{G}\circ\scr{F})^d(a_d,\ldots,a_1)
\\
&\hspace{2em}=\sum_r\sum_{s_1,\ldots,s_r}\scr{G}^r
\bigl(\scr{F}^{s_r}(a_d,\ldots,a_{d-s_r+1}),\ldots,\scr{F}^{s_1}(a_{s_1},\ldots,a_1)\bigr),
\end{split}
\end{equation}
其中的求和如\eqref{eq:afunctor-equations} 中那样。这个\trComposition{}是严格结合的，并且
恒等\trFunctor{}（$\scr{F}X=X$，$\scr{F}^1=\id$ 作为 $\hom_{\scr{A}}(X,X)$ 的%
\trEndomorphism，更高阶的项全为零）是一个中性元。
